\section{A preliminary construction with fresh labels}\label{sec:quadratic-route}
The first construction exhibits all the algebraic steps with a
generous label budget. In characteristic two, it repairs the common
identity $zwz=z$ at tuple length $b$ by allocating $b$ new retained
labels. The cumulative cost is quadratic in $r$; the next section
improves this cost by using the residual's special shape on one pool.

\subsection{Fresh-label construction}\label{sec:quadratic}
We now prove that a support-$r$ witness exists in characteristic two when
\[n\ge T(r):=1+2+\cdots+r=\frac{r(r+1)}2,\qquad r<q.\]
This proof is independent of the rook theorem used later. Its ingredients are finite algebra, a position-compatible auxiliary inverse, and an explicit way to simulate that inverse using fresh labels.

Choose the generic coefficients in a finite field $F$ of characteristic two, large enough to contain $a$ independent elements. We may work inside the algebraic closure of the original field; the degree question descends by Lemma~\ref{lem:field}. The finiteness of $F$ will be used only to normalise an element of the finite unital group algebra on the current label set.
If that label set is $\Omega$, the algebra $F[S_\Omega]$ has
$|F|^{|S_\Omega|}$ elements, one coefficient in $F$ for each
permutation. This finite construction field imposes no finiteness
restriction on the original coefficient field.

\subsubsection{Normalisation of the residual}
Suppose $w_0$ works on shorter tuples but not yet on length $b$. Put $y=z-zw_0z$. An auxiliary linear map $h$ may satisfy $zhz=z$ on the larger tuple space. It is tempting to conclude $yhy=y$. In general the expansion contains an extra term $zw_0zw_0z$, and this need not collapse. Before correcting the error we arrange the additional identity $w_0zw_0=w_0$ in the actual group algebra. This normalisation was a genuine mathematical step in the earlier construction.

Here normalisation replaces a ring element to impose an additional
algebraic identity. A ring representation interprets ring elements as operators,
respecting addition, multiplication and the identity. Thus a relation
proved in the group algebra holds simultaneously in all its
representations. This is why the next lemma preserves every previously
successful tuple length at once.

\begin{lemma}[Normalisation in a finite ring]\label{lem:normal}
Let $a,z$ belong to a finite ring. There is an element $b$ in the same ring with $bzb=b$, such that every ring representation in which $zaz=z$ also satisfies $zbz=z$.
\end{lemma}
\begin{proof}
Put $x=az$. The sequence of powers of $x$ eventually repeats, so $x^{i+t}=x^i$ for some $i\ge0,t>0$. Multiplication by further powers gives $x^{m+t}=x^m$ for $m\ge i$. Choose $k\ge i$ with $k+1$ divisible by $t$. Then $2k+1$ differs from $k$ by a multiple of $t$, and $x^{2k+1}=x^k$.

Set $b=x^ka$. Direct multiplication gives
\[bzb=x^k(az)x^ka=x^{2k+1}a=x^ka=b.\]
In a representation where $zaz=z$, induction gives $z(az)^m=z$ for every $m\ge0$. Hence $zbz=z(az)^{k+1}=z$. One exponent chosen in the original finite ring works simultaneously in all such representations.
\end{proof}
The notation $b$ in this lemma is an element, not a tuple length; in the construction below we denote this normalised element by $b_0$.

\begin{lemma}[Normalised residual]\label{lem:residual}
If $b_0zb_0=b_0$ and $zhz=z$, then, for $y=z-zb_0z$,
\[yhy=y.\]
\end{lemma}
\begin{proof}
Expand without commuting any factors:
\begin{align*}
yhy&=zhz-zhzb_0z-zb_0zhz+zb_0zhzb_0z\\
&=z-2zb_0z+zb_0zb_0z\\
&=z-zb_0z=y.
\end{align*}
The last equality uses the normalisation. This identity works in any ring; characteristic two is not needed for this particular calculation.
\end{proof}

\subsubsection{An auxiliary inverse compatible with tuple positions}
Position permutations reorder the entries of a tuple; label
permutations rename those entries. These actions commute. In this
step we require compatibility with positions, since the subsequent
simulation must work on all orderings of each underlying label set.

Let $\Omega$ consist of the block $B$ and the retained labels already allocated. On $Q=P_b(\Omega)$, partition the tuples according to the $A$-orbit of their \emph{underlying label set}. A set meeting $B$ has proper stabiliser because it has size $b<q$. Choose the functional $\ell$ and coordinate $\pi$ as in Section~\ref{sec:block2}, now on these sets. Let $D$ read $\pi$ from the underlying set, and use the same $W$ and nonzero $\alpha$ as before, defined separately on each orbit. Then on each moved orbit,
\[Dz-zD=W,\quad Wz=zW,\quad W^2=\alpha^2I,\quad z^2=0.\]
Multiplying the commutator by $z$ shows $zDz=zW$. Thus
\[h=\alpha^{-2}WD\quad\text{satisfies}\quad zhz=z.\]
On tuples wholly outside $B$, put $h=0$, since $z=0$ there. Because $D$ depends on the underlying set and translations act on labels, $h$ commutes with every permutation of tuple positions. We have not claimed that $h$ is a group-algebra element. It is just an explicitly constructed linear operator.

\subsubsection{Fresh labels simulate any position-compatible operator}
Adjoin $b$ new labels $\kappa_1,\ldots,\kappa_b$ to $\Omega$. Let $i:P_b(\Omega)\to P_b(\Omega\cup C)$ include the old tuples and let $p$ discard every tuple using a fresh label. We prove that any operator $h$ commuting with position permutations has the form
\begin{equation}\label{eq:compression}h=pci\end{equation}
for some $c\in F[S_{\Omega\cup C}]$.

The inclusion and projection satisfy $pi=I$. The expression $pci$
is the old-to-old part of the larger operator $c$: it makes no claim
about inputs or outputs containing fresh labels. The residual
localisation in the induction below is what makes this part sufficient.

Fix a $b$-element set $U=\{u_1,\ldots,u_b\}\subseteq\Omega$, with one chosen ordering. Form
\[E_U=\prod_{j=1}^b\bigl(1-(u_j\ \kappa_j)\bigr).\]
The transpositions are disjoint. On an old tuple missing some $u_j$, the corresponding factor kills the tuple, and the other factors cannot introduce $u_j$. On a tuple using all the $u_j$, the identity term leaves the tuple old, while every nonidentity term introduces a fresh label. Therefore $pE_Ui$ is precisely the projector onto the orderings of $U$.

For an ordered output tuple $v$, choose a permutation $g\in S_\Omega$ taking $(u_1,\ldots,u_b)$ to $v$. Extend $g$ by fixing fresh labels. Then $pgE_Ui$ takes that input ordering to $v$, takes its position permutations to the corresponding permutations of $v$, and kills all other underlying sets. The orderings of each $U$ form a free orbit for $S_b$. Consequently an operator commuting with $S_b$ is completely determined by its value on one ordering for each $U$. Expanding those values as linear combinations of output tuples and using the operators $pgE_Ui$ proves \eqref{eq:compression}. There is no division by $b!$.

A free position orbit means that no nonidentity position permutation
fixes an injective tuple. Its $b!$ orderings are distinct basis vectors.
Consequently position compatibility determines the remaining values
from one ordering without averaging over $S_b$.

\begin{example}[Two registers]
For $U=\{u_1,u_2\}$,
\[E_U=1-(u_1\kappa_1)-(u_2\kappa_2)+(u_1\kappa_1)(u_2\kappa_2).
\]
On $(u_1,u_2)$, projection $p$ kills the last three terms. On $(u_1,v)$ with $v\ne u_2$, the identity contribution is cancelled by the second swap factor, which fixes that tuple. The detector recognises the set $U$, while leaving the two possible orderings available for a position-compatible output rule.
\end{example}

\subsubsection{The induction, including localisation of the error}
For support zero take $w_0=0$: $z$ is zero on constants. Suppose a witness exists on the block and $T(b-1)$ retained labels, collectively $\Omega$. Normalise it in the finite ring $F[S_\Omega]$ to $b_0$. It still works on all tuple lengths below $b$ and satisfies $b_0zb_0=b_0$ in the ring itself. Set $y=z-zb_0z$.

Adjoin $b$ fresh labels. Decompose $P_b(\Omega\cup C)$ according to the positions occupied by fresh labels and their actual values. The group $S_\Omega$ fixes each fresh label, so every summand with at least one fresh label is a copy of $P_j(\Omega)$ for $j<b$. On it $y=0$ by the old witness relation. The summand with no fresh labels is $Q=P_b(\Omega)$. Thus, as actual maps,
\[y=i y_Q p.\]
This localisation is what allows a compressed inverse to suffice. Construct the auxiliary $h$ on $Q$ above. Lemma~\ref{lem:residual} gives $y_Qhy_Q=y_Q$. Choose $c$ with $pci=h$. Then
\[ycy=i y_Q(pci)y_Qp=i y_Qp=y.\]
Define the new group-algebra element
\[w_b=b_0+(1-b_0z)c(1-zb_0).\]
Since $z(1-b_0z)=(1-zb_0)z=y$,
\[zw_bz=zb_0z+ycy=zb_0z+y=z\]
on length $b$. On every shorter tuple space the same fresh-label decomposition makes $y=0$, so the relation remains valid there too. The new cost is exactly $b$ retained labels, giving
\[T(b)=T(b-1)+b.\]
If still more retained labels are present, extend the constructed permutations by fixing them. Decomposing according to those outside labels reduces every tuple of length at most $b$ to a shorter tuple on the constructed labels. Thus the same witness works for all $n\ge T(b)$.

\begin{theorem}[Quadratic retained-label budget]
For every $r<q=2^a$ and $n\ge r(r+1)/2$, there is a common witness $w$ with $zwz=z$ on all tuple modules of length at most $r$. Therefore deletion is exact on the support-$r$ class. It transports the
entire degree-$d$ PC consequence space with $r=kd+s$, beyond the fixed
template threshold. For prescribed indexed targets use the budget
$r=2kd+s+t$ to include the target map and the doubled-support detector.
\end{theorem}
Every construction needed for this theorem has now been supplied. In particular, the proof is not an inference from witnesses found for small tuple lengths.


\begin{remark}[The square-root degree lower bound]
For either prescribed-target budget, the construction gives an onset
$O(d^2)$ for degree-$d$ transfer.
Inverting that onset gives an $\Omega(\sqrt n)$ lower bound when no
finite degree recurs, and on the hard residues of the periodic case.
The loss is precisely the cumulative allocation $T(r)=r(r+1)/2$.
The following construction replaces this sum by the maximum $2r-1$,
using one retained pool at every layer.
\end{remark}
